% !TEX root = ../math.tex
\chapter{Render Trie}
\label{chap:render-trie}

\section{Scrolling}
This implementation of Dotter is designed for typing 
 medium-length strings of text of about 400 characters or less.
It is intended for typing text messages, emails, chat-bot commands,
and search queries. These are typically a few sentences long.

This raises the problem of scrolling, because 400 characters
is far too much to fit horizontally on a single screen.
Scrolling must be automatic in order not to slow down the user.
The advantage of scrolling is that keeps the active part of the trie
centered and maximizes utilization of real estate. The disadvantage
of scrolling is that it is jarring and disorienting. Below we discuss a simple
heuristic to balance this trade-off.

\subsection{A Quadratic Cost Heuristic}

We assume that all characters occupy constant width.
We formalize our goal of keeping the active part of the trie centered
in terms of the ``first fork.'' The first fork node is the first node
that has multiple children in the visible trie. We denote its depth
as $d$. We define the scroll offset $s$ as the number of characters
which are off-screen to the left of the window.

Our centering objective means that we want to 
minimize the distance between the
``first fork position'', $p := w(d - s)$, where $d$ is the depth of the first fork node and $w$ is the width of any node which is a single parent, and the desired first fork position, $p^*$.
Our stability objective means that we want to minimize 
changes in the offset between iterations, $\Delta s = s_t - s_{t-1}$.

We may parameterize such an objective with rates $a$ and $b$,
and solve for $s_t^*$,

\begin{align*}
    J &= a (\frac{p^* - p}{w})^2 + b \Delta s^2 \\
    J &= a (s_t - (d - \frac{p^*}{w}))^2 + b (s_t - s_{t-1})^2 \\
    \implies 0 &= 2a (s_t - (d - \frac{p^*}{w})) + 2b (s_t - s_{t-1}) \\
    \implies s_t^* &= \frac{a (d - \frac{p^*}{w}) + b s_{t-1}}{a + b}
\end{align*}


\subsection{Two Approaches to Backspace}

There are two ways to allow the user to go back to 
nodes that are beyond the scroll window. Either we can have
a single "backspace" node that contributes likelihood to all
off-screen predecessors, or we can try to actually assign independent
timers to every off-screen predecessor. Not only is the former
less cluttered, it also has the advantage of showing incremental progress.

\section{Embellishments}
\subsection{Branches of the Tree}
% curvature is minimized with cubic splines
% the constant velocity condition is that r'(t) \dot r''(t) = 0
% There are eight dof
% start and end positions define 4 of them
% initial and final angles define 2
% equal initial and final speeds define 1
% constant velocity defines 1

% the speed implicity defines the length of the curve
% we can calculate the curvature as 12a_3^2 + 12a_2a_3 + 4a_2^2
% but then we'd like to imagine reparameterizing so that
% the initial speed is unit.
% saying that the particles is now traveling slower by a factor of [speed]
% means that the total squared curvature decreases by a factor of [speed]^2

% N.B. Euler's method produces necessary conditions (arising from δJ/δf = 0)
% but for J to be an extremum this condition is not sufficient
% We would also need to ensure that the "Hessian" of variations is positive definite
